\chapter{整数下降怎样跨过合数障碍}\label{ch:tail}
\section{非单位差分也能帮忙}
前篇常用差分 $a-b$ 的模逆来调配比。现在已有至少 $k$ 份 $a$ 和至少 $k$ 份 $b$，且 $a\ne b$。对合数 $k$，设 $g=\gcd(a-b,k)>1$，模逆不存在，但只取 $k/g$ 份 $a$、其余取 $b$ 就得到
\begin{equation}\label{eq:nonunit}
\frac{(k/g)a+(k-k/g)b}{k}=b+\frac{a-b}{g}\in\Z.
\end{equation}
两种值都留有未选副本，因此这一步安全。\src{S04}

例如六平均中 $a=5,b=1$，差分四与六的最大公因数为二。取三份五、三份一，得到六份三。不可逆的差分没有让我们无路可走，反而给出了一个只用两种值的整数平均。

\section{偶数元数唯一新增的碰撞：两个配比同为中点}
把两种配比对调，两个新均值是
\[c_1=b+(a-b)/g,\qquad c_2=a+(b-a)/g.
\]
其差为 $(2/g-1)(a-b)$，所以只在 $g=2$ 时相等，此时都是中点。

若恰有三个重值 $a<b<c$，原本想用两种配比之一避开第三值。唯一可能同时撞上第三值的情形，是两端平均到中间：
\[a+c=2b,\qquad\gcd(c-a,k)=2.
\]
此时相邻差分必须是单位；若相邻差分也非单位，其两值混合的均值严格在相邻两值之间，不会等于第三个。于是又可以用前篇的大小顺序论证：有轻值就保留另一极端，极端有富余就只取它一份；最后只剩相反值加零的终端。

\section{其余分类怎样接回前篇}
大背景值的证明分两支：差分非单位时直接用\eqref{eq:nonunit}；差分单位时仍可唯一选择 $i$，使 $(a^i,b^{k-1-i},x)$ 为整数平均。若所有选择碰撞，所有值被迫在同一有限等差列中，前篇的系数跨越打包继续成立。

需要补的一点是，等差列中的 $x-b=-j(a-b)$ 也可能是非单位。遇到它不能继续当作可逆差分；再次使用\eqref{eq:nonunit}即可。三个以上重值用上一节的中点分析，恰两个重值则各留 $k$ 份，再保护 $d_k(n)$ 个见证，在剩余池应用一般EGZ。

\begin{lemma}[一般EGZ由素数版本推出]
任意 $2k-1$ 个整数中，有 $k$ 项的和被 $k$ 整除。
\end{lemma}
\begin{proof}
素数情形已在前篇证明。若 $k=uv$，从 $2uv-1$ 项中反复取出 $2v-1$ 个不交的 $u$ 元组，使每组和被 $u$ 整除；最后一次取组前仍至少有 $2u-1$ 项。将各组和除以 $u$，应用 $v$ 的版本，选 $v$ 组的并集。对因子个数归纳。
\end{proof}

\begin{theorem}[任意元数的已证尾部]
对 $k\ge3$，若
\[n\ge4k-1+d_k(n),\qquad d_k(n)=\#\{\ell\mid n:\ell\nmid k\},
\]
则 $\mathcal P(k,n)$ 成立。特别
\[N(k)\le4k+\lceil\log_2k\rceil+1\le5k.
\]
\end{theorem}
\begin{proof}
从 $kX$ 初始化两个重值。上面已核对前篇分类的全部合数改动。恰两重值且都不足 $2k-1$ 份时，各留 $k$ 份，再保护至多 $d_k(n)$ 个位置，余池至少 $2k-1$ 项，且每个实际值不足 $k$ 份；一般EGZ给非恒等安全整数平均。
因此每个非终端都有保持储备的整数续步，平方和严格下降，有限步进入相反值终端。最后的统一界用 $d_k(n)\le\lfloor\log_2n\rfloor$ 及整数函数 $n-\lfloor\log_2n\rfloor$ 不减，与前篇相同。
\end{proof}

这项结论已经保证每个 $k$ 的最终阈值\emph{存在}。未完成的是精确的最优阈值。旧文曾说合数还不知道是否有最终阈值，那是早于本定理的阶段。

\section{精度与算法：这里究竟保证多快}
对合法整数零和 $X$，在 $kX$ 上做整数下降，原 $X$ 的中间值都属于 $(1/k)\Z$。合数情况下，分母可以是 $k$ 的真因子，例如六平均可出现二分之一、三分之一；不能沿用“分母只有一或 $p$”的素数断言。

整数能量每次至少降二，给粗界 $T\le(k^2/2)\sum X_i^2$。这是关于数值大小的界，不是输入位长的多项式界。已知充分范围内的\emph{判定}则很快：算 $G$，反复除去 $\gcd(G,k)$，直到得到一或剩余与 $k$ 互素；不需要分解巨大 $G$。

四平均已有 $N(4)\le16$，六平均已有 $N(6)\le24$。还有一个方便终端：对这两种元数，任意维数一旦产生 $k$ 个零就能完成。非零支持不超过 $k$ 时直接一次平均；其他情况四平均保留一个零后用偶数维数转译，六平均保留 $n\bmod3$ 个零后用三倍数维数定理。含零子系统自身零和且本原 $G=1$，调用合法。\src{S04}
